\chapter{Flux Reconstruction}\label{ch:11}
\begin{chaptersummary}
The strong DG form represents interface residuals as liftings into the element interior. FR follows this idea by interpolating the local flux and incorporating a shared flux through endpoint correction functions. This chapter derives the nodal update and conservation property, then identifies the exact relationship with a particular DG discretization through a mass-matrix lifting condition.
\end{chaptersummary}
\learninggoals{Construct discontinuous and corrected fluxes for a one-dimensional conservation law; write the FR spatial operator; prove cell-average conservation; state the basis, integration, and flux conditions required for a DG equivalence.}
\studyroute{Multidimensional curved meshes, nonlinear aliasing, and general FR entropy analysis are omitted. The next chapter constructs energy-stable correction functions.}{Huynh flux reconstruction; correction function; lifting operator; SBP correction procedure}{At three Gauss nodes, construct a quadratic solution and cubic correction functions, assemble the nodal FR operator, and compare it with exactly integrated strong-form DG.}

\section{Discontinuous Fluxes and Endpoint Corrections}\index{FR methods}
Consider $u_t+\partial_xF(u)=0$ on an affine element $x=x_j+J_j\xi$.
Choose $p+1$ distinct nodes $\xi_i$ with Lagrange basis $\ell_i$:
\[
u_j^\delta(\xi)=\sum_{i=0}^pu_{ji}\ell_i(\xi),\qquad
f_j^D(\xi)=\sum_{i=0}^pF(u_{ji})\ell_i(\xi).
\]
Here $f^D$ is the degree-$p$ flux interpolant; the nonlinear flux $F(u^\delta)$ may have higher degree.
An interface solver uses the two solution traces to produce shared, consistently oriented fluxes $f_L^*,f_R^*$.
Choose correction functions of degree at most $p+1$ satisfying
\[
g_L(-1)=1,\quad g_L(1)=0,\qquad
g_R(-1)=0,\quad g_R(1)=1.
\]
The corrected flux is
\begin{equation}\label{eq:fr-flux}
f^C=f^D+(f_L^*-f^D(-1))g_L+(f_R^*-f^D(1))g_R.
\end{equation}
Thus $f^C(-1)=f_L^*$ and $f^C(1)=f_R^*$.
The globally shared interface flux has the same endpoint value in each local polynomial. For the original FR construction, see~\cite{huynh2007}.

\conceptbridge{How can two endpoint flux values change every interior degree of freedom?}
{Correction functions extend endpoint flux differences into element-interior polynomials. Their derivatives deliver interface residuals to every solution node. The correction shape determines how this lifting acts on polynomial modes and therefore affects the spectrum and stability.}
{Endpoint flux differences $\longrightarrow$ correction polynomials $\longrightarrow$ interior residual lifting.}

\section{The Nodal FR Update and Local Conservation}
Let $D_{ik}=\ell_k'(\xi_i)$, $C_{iL}=g_L'(\xi_i)$, and $C_{iR}=g_R'(\xi_i)$.
Constrain the corrected differential residual at the nodes:
\begin{equation}\label{eq:fr-operator}
J_j\dot u_{ji}
=-\sum_kD_{ik}F(u_{jk})
-(f_L^*-f_L^D)g_L'(\xi_i)
-(f_R^*-f_R^D)g_R'(\xi_i).
\end{equation}
Let $R$ be the $2\times(p+1)$ endpoint-evaluation matrix. Then
$J_j\dot{\bm u}_j=-D\bm f_j-C(\bm f^*-R\bm f_j)$.
All derivatives have degree at most $p$, so the nodal equality is also a polynomial identity throughout the element.
Exact integration gives
\[
\frac{\dd}{\dd t}\left(J_j\int_{-1}^1u_j^\delta\dd\xi\right)
=-f_R^*+f_L^*.
\]
Shared fluxes and endpoint conditions therefore ensure cell-average conservation. The particular correction functions must additionally satisfy the desired stability conditions.

\section{The Matrix Relationship with DG}
In the same basis, let $M_{ik}=\int\ell_i\ell_k$ and $B=\diag(-1,1)$.
Polynomial integration by parts gives
\[
MD+D^{\mathsf T}M=R^{\mathsf T}BR.
\]
For a linear flux $F(u)=au$, strong-form DG is
\[
J_j\dot{\bm u}_j=-aD\bm u_j-M^{-1}R^{\mathsf T}B
(\bm f^*-aR\bm u_j).
\]
Therefore, when
\begin{equation}\label{eq:fr-dg-lift}
C=M^{-1}R^{\mathsf T}B,
\end{equation}
and representation, integration, mapping, and fluxes agree, FR and this semidiscrete DG operator coincide.
An invertible basis transformation gives the same relationship in a modal implementation.
Nonlinear fluxes, underintegration, or other correction functions require a new comparison of volume and boundary operators.

\section{Error and Spectral Analysis}
Correction functions change the interface feedback into element modes. On a uniform periodic mesh, the Bloch symbol gives the spectrum,
allowing comparison of dissipation, dispersion, and explicit stable steps on the physical branch.
An energy method instead examines evolution directly in a chosen norm and includes physical boundaries.
Dissipation, dispersion, and stable step size also depend on fluxes, correction functions, and the time integrator. Cost comparisons by total degrees of freedom and error therefore record these choices alongside formal order.

\aiexercise{For $p=1$, a linear flux, and Gauss solution nodes, calculate $M,D,R,B$. Obtain $g_L',g_R'$ from~\eqref{eq:fr-dg-lift}, integrate them, check all four endpoint conditions, and compare the local matrices with DG.}
{Replace a correction function by another polynomial satisfying the same endpoint conditions. Check conservation, then calculate the periodic-symbol eigenvalues to distinguish the requirements for conservation and stability.}
